% Conclusion.
\section{Conclusion}
\label{sec:conclusion}

We showed that the higher-order geometry of the reconstructed phase-space
attractor is precisely what transfers across a motor-load shift: on the CWRU
single- and multi-source protocols the structural models and GA-DCHG sit
\GapCZeroThreeDynhg{}/\GapCZeroThreeGad{} and
\GapCZeroOneTwoThreeDynhg{}/\GapCZeroOneTwoThreeGad{} points above a raw 1D-CNN,
which degrades sharply under the shift. Building on this, GA-DCHG uses the
edge-geometry descriptors as anchors rather than passive gates. The anchors
screen a per-protocol anchor, modulate hyperedge generation, align source and
target in geometry space, and validate decisions by intervention auditing.
GA-DCHG attains the best target-test accuracy of all models on both protocols.
The per-protocol anchor screening is itself a measurable result, the component
ablation on the multi-source protocol confirms the anchor and temporal components
carry part of the cross-domain benefit, and the intervention audit evidences
that the cross-domain decision is sensitive to the edge geometry. We release
all code so that every reported number is reproducible from the raw recordings.
